% Quantitative component bounds used by the dual wire-symbol estimate.
\section{Component estimates for the dual wire symbol}
\label{sec:ex-lattice-free-kernel-components}

This section supplies the factor bounds used in the dual-stencil calculation
of Section~\ref{sec:ex-lattice-free-kernel-estimates}. All constants below
may depend on the fixed smoothing stages, the positive coupling parameter,
and the four-dimensional cell conventions, but not on the odd side length
or the alias. Write
\[
 \Gamma_L(l)=\prod_{a=1}^4(1+|\ell_L(l)_a|),\qquad
 g_L(l)=1+\max_a|\ell_L(l)_a|.
\]
The centered representative satisfies $g_L(l)\leq\Gamma_L(l)$. For
$z\in\mathbb C$ define
\[
 \operatorname{rail}_L(z)=\frac1L\sum_{v=0}^{L-1}e^{ivz},\qquad
 \operatorname{plank}_L(z)=e^{-i\lfloor L/2\rfloor z}
                            \operatorname{rail}_L(z).
\]
For $p\in\mathbb C^4$ set
\[
 \operatorname{area}_L(p)=\prod_{a=1}^4\operatorname{plank}_L(p_a).
\]
Set $\mathrm{zone}=\pi/4$ and $R_0=3\pi/2$. Since
$(e^{iz}-1)/z$ extends continuously and has no
zero on the closed disk $|z|\leq R_0$, compactness gives constants
$0<\ell\leq h$ such that
\[
 \ell|z|\leq|e^{iz}-1|\leq h|z|\qquad(|z|\leq R_0).
\]
Put $p_0=\min(1,\ell/h)$ and $c_0=hR_0/p_0$. These are the uniform
rail/plank lower and principal-plank constants used below.
In the divisor estimates write $\mathrm{low}=\ell$ and $\mathrm{high}=h$;
for real four-vectors, $|p|$ denotes the Euclidean norm.
At each stage $j$, a train is built by finite recursion. The top train
$\mathrm{top}_j$ has mass side $1$ and extra factor $1$. A hop has output
stage $j$, coarse dilation an odd integer $L\geq3$, and a predecessor train
of stage $J\in\{j,j+1\}$; if the predecessor has mass side $M$, the hop has
mass side $LM$ and extra factor
$\operatorname{area}_{LM}(p)^{J-j}$ times the predecessor's extra factor.
Thus every train is a finite history, and no bound on its number of hops is
assumed.
For a side $M$, set
\[
 s_M(p)_a=M\bigl(e^{ip_a/M}-1\bigr),\qquad
 d_M(p)=\sum_{a=1}^4s_M(p)_a s_M(-p)_a.
\]
Thus $d_M$ is the divisor, and for real $p$ it is the sum of the four
squared seed magnitudes. A train $v$ of stage $j$ and mass side $M=v.mass$
has cumulative filter $S_v(p)=\operatorname{area}_{M}(p)^jv.extra(p)$; this is distinct
from the patch mass factor $m_P(p)=\operatorname{area}_{P.L}(p)^{P.J}$
used in the wire branches. A flight with indices $j,t\geq4$ consists of an
odd coarse side $L\geq3$, a wing train of stage $j$, and indices
$t\leq j\leq t+1$; its land train is the hop of that wing with output
stage $t$. Its patch mass is the separate factor
$m_P=\operatorname{area}_L^{\,j}$.
For a train of side $M$ and direction $a$, define its drive, lid, and halo by
\[
 D_{v,a}(q)=S_v(q)\operatorname{rail}_M(q_a),\quad
 \ell_{v,a}(p)=D_{v,a}(p/M)D_{v,a}(-p/M),\quad
 h_v(p)=S_v(p/M)S_v(-p/M).
\]
For an alias side $L$, write $q_l^{(L)}(p)=(p+2\pi\ell_L(l))/L$. For a train of side $M$, set $\operatorname{blow}_{M,L}(p,l)=L^2d_M(q_l^{(L)}(p))$. The lift bound $|\ell_L(l)_a|\leq(L-1)/2$ gives $q_l^{(L)}(\mathcal Q_\eta)\subset\mathcal Q_\eta$. The feed with constant
input one, the high-alias term, the noise drain, bulb, and bend are
\[
\begin{aligned}
 F_{v,a}(p)&=M^{-2}\sum_{l\in\Lambda_M}
       D_{v,a}(q_l^{(M)}(p))D_{v,a}(-q_l^{(M)}(p)),\\
 H_{v,a}(p)&=\sum_{l\ne0}
       \frac{D_{v,a}(q_l^{(M)}(p))D_{v,a}(-q_l^{(M)}(p))}
            {M^2d_1(q_l^{(M)}(p))},\\
 G_{v,a}(p)&=H_{v,a}(p)+k^{-1}N_{v,a}(p),\\
 b_{v,a}(p)&=\ell_{v,a}(p)+d_M(p)G_{v,a}(p),\\
 \beta_v(p)&=h_v(p)^{-1}-\sum_{a=1}^4
      \frac{J_a(p)G_{v,a}(p)}{\ell_{v,a}(p)b_{v,a}(p)},
 \qquad J_a(p)=(e^{ip_a}-1)(e^{-ip_a}-1).
\end{aligned}
\]
The noise is defined from a list: the top train has empty list; a hop of
coarse side $L$ has as its list one top train of side one followed by the
images of every entry of the predecessor's list under that same hop.
For a list $\mathcal N$, its noise at direction $a$ is the sum, over
$u\in\mathcal N$, of the feed
\begin{align*}
 \operatorname{feed}_{u,a}(p)&=\frac1{u.mass^2}
   \sum_{l\in\Lambda_{u.mass}}
      D_{u,a}(q_l^{(u.mass)}(p))D_{u,a}(-q_l^{(u.mass)}(p)),\\
 N_{v,a}(p)&=\sum_{u\in\mathcal N(v)}\operatorname{feed}_{u,a}(p).
\end{align*}
In the high-alias term, the divisor side
is the unit side, as in the source definition of the high part of the
drain; the port blow-up below uses the indicated train side $M$.

\subsection{The one-dimensional alias estimate}

\begin{lemma}[Rail and box decay]
\label{lem:ex-lattice-rail-alias-decay}
There is a constant $B$ such that, for every odd $L\geq1$, every
$p\in\mathcal Q_\eta$ with $\eta$ in a fixed sufficiently small cube, and
every centered residue $l$,
\[
 |\operatorname{rail}_L(q_l^{(L)}(p)_a)|,
   |\operatorname{plank}_L(q_l^{(L)}(p)_a)|
 \leq B(1+|\ell_L(l)_a|)^{-1}.
\]
At the zero residue the reciprocal rail is also uniformly bounded on a
common smaller complex cube.
\end{lemma}

\begin{proof}
If $\ell_L(l)_a=0$, the imaginary part of $Lq_l^{(L)}(p)_a=p_a$ is bounded by
$\eta$, and the normalized finite geometric sum defining the rail has
modulus at most $\exp(\eta)$; multiplication by the centered phase defining
the plank costs at most another fixed exponential factor. If
$\ell_L(l)_a\ne0$, the centered-lift estimate for aliases gives
\[
 (\pi/4)(1+|\ell_L(l)_a|)\leq
 L|q_l^{(L)}(p)_a|\leq C(1+|\ell_L(l)_a|).
\]
The lower bound uses that $|\operatorname{Re}p_a|<\pi+\eta$ and the lift is
a nonzero integer; the upper bound is the triangle inequality. On the fixed
disk containing these momenta, the function
$(e^{iz}-1)/z$ is bounded above and bounded away from zero, with its value at
zero defined by continuity. The identity
\[
 L\,\operatorname{rail}_L(z)(e^{iz}-1)=e^{iLz}-1
\]
then bounds the rail by a constant divided by
$L|e^{iz}-1|$, hence by the stated inverse lift weight. The plank differs
from the rail by a centered phase whose modulus is bounded on the same
complex cube. At the zero alias, the real-frequency rail has a uniform
positive lower bound on the contracted disk. The following quantitative
continuation promotes it to a common complex cube, uniformly in the side.
Suppose a holomorphic function has modulus at most $Cw$ on a cube of width
$\eta$, with $w>0$, and modulus at least $cw$ on its real slice of width
$\theta$. At a complex point $p$, let $y=\operatorname{Re}p$. The
one-variable Schwarz estimate on a ball centered at $y$ gives
\[
  |f(p)-f(y)|\leq \frac{4Cw}{\eta}\,\|p-y\|_\infty
\]
provided $\|p-y\|_\infty\leq\eta/4$ and the ball lies in the original cube. Thus on
the smaller cube of width
$r=\min(\eta/4,\theta,c\eta/(8C))$ one has
$|f(p)|\geq(c/2)w$. In particular the reciprocal is holomorphic there and
bounded by $2/(cw)$. Applying this with $w=1$ to the zero-alias rail proves
the asserted uniform reciprocal bound.
\end{proof}

\begin{lemma}[Principal plank estimate with unit coefficient]
\label{lem:ex-lattice-principal-plank}
If $M\leq L$, $z$ is real, and $L|z|\leq R_0$, then
\[
 |\operatorname{plank}_M(z)|\geq
 \exp\bigl(-c_0 M/L\bigr).
\]
\end{lemma}

\begin{proof}
The centered-sum formula is
\[
 \operatorname{plank}_M(z)=\frac1M\sum_{v=0}^{M-1}
   e^{i(v-\lfloor M/2\rfloor)z}.
\]
Each centered integer has modulus at most $M$. The chip upper bound
therefore gives
$|\operatorname{plank}_M(z)-1|\leq hM|z|\leq hR_0x$, where $x=M/L$.
The rail-chip quotient and the two-sided chip comparison also give
$|\operatorname{plank}_M(z)|\geq p_0$ because $M|z|\leq R_0$ and
$|z|\leq R_0$. Consequently
$1\leq|\operatorname{plank}_M(z)|+hR_0x$ by the triangle
inequality. Since $e^{c_0x}\geq1+c_0x$ and
$|\operatorname{plank}_M(z)|/p_0\geq1$, it follows that
$e^{c_0x}|\operatorname{plank}_M(z)|\geq
|\operatorname{plank}_M(z)|+hR_0x\geq1$.
Rearranging proves the claimed estimate. The coefficient $c_0$ is fixed
independently of the number of train steps; this is why its exponential
loss can be summed by the geometric side-ratio argument.
\end{proof}

\subsection{Mass, signals, and nonzero-alias denominators}

\begin{lemma}[Mass and flux weights]
\label{lem:ex-lattice-mass-flux-weight}
For a flight $P$ with patch side $L$ and stage $j$, its patch mass factor
$m_P(q)=\operatorname{area}_{L}(q)^j$ obeys, on a common complex cube,
\[
 |m_P(q_l^{(L)}(p))|\leq C_m^j\Gamma_L(l)^{-j}.
\]
The $r$-face flux has entrywise bound
$C_{m,r}\Gamma_L(l)^{-j}$ for each fixed $r$ used in the wire formula.
At the zero alias, the latch factor $m_P(q_0)^{-1}$ and reciprocal rails
are bounded independently of $L$ on a common smaller complex cube.
\end{lemma}

\begin{proof}
The area is the product of the four one-dimensional plank factors. Apply
Lemma~\ref{lem:ex-lattice-rail-alias-decay} in each coordinate and raise
the result to the $j$th power. The product of the four coordinate weights
is exactly $\Gamma_L(l)^{-j}$. A flux entry is this mass evaluated at
$-q_l$ times the corresponding rail factors at $-q_l$, followed by
transpose; the rail factors have uniformly bounded modulus, so the same
weight remains. At zero alias, the patch mass obeys
\[
 |m_P(q_0)|\geq \exp(-4j\,c_0)
\]
on the real cube: apply the principal one-dimensional plank lower bound
with scale ratio one in each coordinate. Its upper analytic-cube bound
and this real lower bound give a uniform reciprocal on a smaller complex
cube by the quantitative continuation estimate above. The zero-alias
rails are handled by their real lower bound and the same continuation.
\end{proof}

\begin{lemma}[Scaled signal weight]
\label{lem:ex-lattice-scaled-signal-weight}
For every face degree $r$, the scaled signal has entrywise bound
$C_sg_L(l)$, independent of $r$. In particular the one-form signal costs
at most $C_sg_L(l)$.
\end{lemma}

\begin{proof}
The signal is exterior multiplication by the vector with entries
$L(e^{iq_l^{(L)}(p)_a}-1)$. Its matrix entry from an $r$-face to an
$(r+1)$-face is zero unless the latter face is obtained by adjoining one
coordinate; in that case it is, up to sign, the corresponding single vector
entry. The alias upper bound and $|e^{iz}-1|\leq |z|e^{|\operatorname{Im}z|}$
bound this entry by $C(1+|\ell_L(l)_a|)\leq Cg_L(l)$. Thus the cost is one
grade factor for every $r$.
\end{proof}

\begin{lemma}[Cumulative train, noise, bulb, and bend bounds]
\label{lem:ex-lattice-train-component-bounds}
Fix a stage $j\geq4$ and $k>0$. Uniformly over every finite train history
of that stage and its mass side $M$, there is a common complex cube on which
for every $p$ in the cube and every alias $l$, the cumulative filter
$S_v(q_l^{(M)}(p))$ and directional drive $D_{v,a}(q_l^{(M)}(p))$ are bounded
by $C\Gamma_M(l)^{-j}$. At the zero alias the real lower bound
$|S_v(p/M)|\geq c>0$ holds uniformly on the fixed real cube. On a possibly
smaller common complex cube, the noise drain, bulb, bend, and all
reciprocals used in the pulse and vault are bounded. On its real slice there
are uniform positive lower bounds for each directional lid, bulb, and
bend. For each nonzero alias $l$, the directional drive therefore has the
stated bound $C\Gamma_M(l)^{-j}$, evaluated at the aliased momentum.
\end{lemma}

\begin{proof}
For a train $v$ of stage $J$ and side $M$, write
$S_v(p)=\operatorname{area}_M(p)^Jv_{\rm extra}(p)$. A hop with output
stage $j$ and predecessor of stage $J\in\{j,j+1\}$ has side $LM$ and
satisfies
\[
 S_{\rm hop}(p)=\operatorname{area}_L(Mp)^J S_v(p),\qquad
 v_{{\rm extra},{\rm hop}}(p)=
     \operatorname{area}_{LM}(p)^{J-j}v_{\rm extra}(p).
\]
The first identity follows by multiplying the four box averages and using
their dilation composition law. For the upper bound, if $p$ lies in the
strip of width $\mathrm{zone}/L_0$ and $M\leq L_0$, each box factor has
modulus at most $e^{\mathrm{zone}M/L_0}$, hence the fourfold product at most
$e^{4\mathrm{zone}M/L_0}$. Induction over histories proves
$|v_{\rm extra}(p)|\leq e^{8\mathrm{zone}M/L_0}$. At a hop with $J=j$,
the extra is unchanged and the mass increases. If $J=j+1$, the added box
factor contributes exponent $4\mathrm{zone}LM/L_0$; the predecessor's
extra is bounded by exponent $(8/3)\mathrm{zone}LM/L_0$ since $L\geq3$.
Their sum is at most $8\mathrm{zone}LM/L_0$.

For the lower bound, the principal one-dimensional estimate with coefficient
one states that if $M\leq L_0$, $z$ is real, and $L_0|z|\leq
R_0$, then
$|\operatorname{plank}_M(z)|\geq
e^{-c_0M/L_0}$. Four coordinates give the same estimate
with exponent multiplied by four. The same induction proves
$|v_{\rm extra}(p/L_0)|\geq e^{-8c_0M/L_0}$.
For $J=j$, the extra is unchanged and the mass increases. For $J=j+1$,
the new box costs exponent $4c_0LM/L_0$, and the old
extra costs at most $(8/3)c_0LM/L_0$; their sum is bounded
by $8c_0LM/L_0$. This is the geometric mass-ratio estimate
that keeps the bound independent of the history length.

For $M=1$ the only alias is zero and the top train has $S_v=D_{v,a}=1$. For $M>1$, the alias estimate for each plank gives
$|\operatorname{area}_M(q_l^{(M)}(p))|\leq B^4\Gamma_M(l)^{-1}$, while the
upper extra estimate at $L_0=M$ costs at most $e^{8\mathrm{zone}}$.
Consequently, for the fixed output stage $j$,
$|S_v(q_l^{(M)}(p))|\leq (B^4)^j e^{8\mathrm{zone}}\Gamma_M(l)^{-j}$.
At zero alias the lower estimate with $L_0=M$ gives
$|S_v(p/M)|\geq\exp(-4j\,c_0-8\,c_0)$.
The zero-root rail is at least $p_0$ and at most a uniform
constant. Thus every directional drive and its two-sided lid have uniform
upper and positive lower bounds there. At any alias, the train smooth upper
bound is $C\Gamma_M(l)^{-j}$ and the rail is bounded, so the drive has that
same weight.

The divisor $d_M(p)$ satisfies, for real $p$ in the fixed disk,
\[
 \mathrm{low}^2|p|^2\leq d_M(p)\leq\mathrm{high}^2|p|^2.
\]
Indeed each seed magnitude lies between $\mathrm{low}|p_a|$ and
$\mathrm{high}|p_a|$, and the divisor is the sum of their squares. The
single-feed noise term uses the constant function one and equals
\[
 M^{-2}\sum_{l\in\Lambda_M}
  D_{v,a}(q_l^{(M)}(p))D_{v,a}(-q_l^{(M)}(p)).
\]
The alias drive bound gives a summand at most
$C\Gamma_M(l)^{-2j}\leq C\Gamma_M(l)^{-2}$, whose sum is at most $81C$;
the retained factor $M^{-2}$ therefore bounds the feed by $81C M^{-2}$.
The source noise list is empty for a top train and, for a hop of dilation
$L$, consists of one top train followed by the hop applied to every entry
of the previous list. Its total price $T$ obeys
$T_{\rm top}=0$ and $T_{\rm hop}=1+L^{-2}T_{\rm old}$. As $L\geq3$,
induction gives $T\leq 1+2/9<2$. Summing the feed estimates at these actual
inverse-square side prices bounds the full accumulated noise uniformly.

The high-alias drain term is a finite sum of two-sided drive products times
the reciprocal blow-up $[M^2d_1(q_l^{(M)}(p))]^{-1}$. For nonzero $l$, choose a
coordinate attaining $\max_a|\ell_M(l)_a|$. The centered-alias lower bound
in that coordinate and the seed lower bound give
$M^2d_1(q_l^{(M)}(p))\geq c\,g_M(l)^2$ on the real cube. Its complex upper bound
has the same $g_M(l)^2$ weight. The quantitative continuation argument
above therefore bounds its reciprocal by $C g_M(l)^{-2}$ on a common
smaller cube. At $M=1$ the high-alias sum is empty. For $M>1$, the remaining drive product has weight
$\Gamma_M(l)^{-2j}$, so this sum is bounded by the exponent-two residue
sum. Thus the full drain is uniformly bounded. It is nonnegative on the
real slice.

The bulb is the directional lid plus $d_M(p)$ times the drain. The real
lower bounds for the lid and nonnegativity of both latter factors give a
positive lower bound for the bulb; the complex upper bound follows from
the already established bounds. The halo is
$S_v(p/M)S_v(-p/M)$ and has a positive real lower bound by the same
cumulative-filter estimate. The exact bend identity follows as follows: the rail-chip relation gives
$J_a(p)h_v(p)=s_M(p)_a s_M(-p)_a\ell_{v,a}(p)$, where
$h_v=S_v(p/M)S_v(-p/M)$ and $\ell_{v,a}$ is the directional lid.
Therefore $d_M(p)/h_v(p)=\sum_aJ_a(p)/\ell_{v,a}(p)$. Combining this
with $b_{v,a}=\ell_{v,a}+d_MG_{v,a}$ yields
\[
 d_M(p)\,\beta_v(p)=\sum_{a=1}^4\frac{J_a(p)}{b_v(p)_a},\qquad
 J_a(p)=(e^{ip_a}-1)(e^{-ip_a}-1).
\]
Indeed, each summand on the right is
$J_a/\ell_{v,a}-d_MJ_aG_{v,a}/(\ell_{v,a}b_{v,a})$, which is exactly
the corresponding term in $d_M\beta_v$. On the real cube, $J_a\geq0$,
$\sum_aJ_a\geq\mathrm{low}^2|p|^2$, and $d_M(p)\leq
\mathrm{high}^2|p|^2$. If $|b_v(p)_a|\leq C_b$, the identity gives
$\beta_v(p)\geq\mathrm{low}^2/(C_b\mathrm{high}^2)$ for $p\ne0$;
at zero $\beta_v(0)=1$. The bend's complex-cube upper bound follows from
its defining formula: first shrink the cube using the positive real lower bounds for the halo and
lid, so their reciprocals are bounded. Next the bulb has a positive real
lower bound and a complex upper bound, so shrink again to bound its reciprocal. With these factors, the displayed bend formula gives a complex
upper bound. The real identity then gives a positive lower bound for the
bend; shrink once more to bound its reciprocal for the pulse and vault. This proves all assertions uniformly over train histories.
\end{proof}

\begin{lemma}[Port denominator and levy weights]
\label{lem:ex-lattice-port-denominator-weights}
For every nonzero alias, the port vector has entrywise weight
$C_Dg_L(l)^{-2}$, the reciprocal port normalizer has weight
$C_Rg_L(l)^{-4}$, and the levy scalar is uniformly bounded when
$j\geq4$.
\end{lemma}

\begin{proof}
The bulb and reciprocal bend bounds are those of
Lemma~\ref{lem:ex-lattice-train-component-bounds}. The blow-up is
$L^2d_M(q_l^{(L)}(p))$, where $M$ is the wing train's side. Choose a coordinate $a_*$ attaining
$\max_a|\ell_L(l)_a|$; since $l\ne0$, its lift is nonzero. The centered
alias lower bound gives
$L|q_l^{(L)}(p)_{a_*}|\geq c(1+|\ell_L(l)_{a_*}|)=cg_L(l)$
on the real cube. The seed lower bound
$|M\,\operatorname{chip}(q_l^{(L)}(p)_{a_*}/M)|\geq
\mathrm{low}|q_l^{(L)}(p)_{a_*}|$ then implies
$L^2d_M(q_l^{(L)}(p))\geq c' g_L(l)^2$. The analytic-cube upper bound for the
blow-up has the same weight $g_L(l)^2$, so the quantitative continuation
argument gives its reciprocal at cost $Cg_L(l)^{-2}$ on a common smaller
cube. The port vector is the bounded bulb times this reciprocal. The
reciprocal port normalizer is the reciprocal of blow-up squared times the
bend; it therefore costs $Cg_L(l)^{-4}$. Finally, the levy contains two
patch mass factors, of total weight $\Gamma_L(l)^{-2j}$, times the bounded
land bend and this reciprocal port normalizer. Since both weights are at
most one, the levy is uniformly bounded.
\end{proof}

\subsection{Uniform tails and regular zero-alias factors}

\begin{lemma}[Uniform tail and regular factors]
\label{lem:ex-lattice-tail-notch-weights}
The tail vector, the correction matrix, the notch matrix, and every
zero-alias dual multiplier in the dual-stencil formula are bounded on a
common complex cube, uniformly in $L$ and the train. The same is true of
the bulb and bend factors and their reciprocals wherever they occur in the
wire formula.
\end{lemma}

\begin{proof}
The train, bulb, bend, and reciprocal bounds are
Lemma~\ref{lem:ex-lattice-train-component-bounds}; the nonzero port bounds
are Lemma~\ref{lem:ex-lattice-port-denominator-weights}. For each nonzero
alias, a tail summand is the product of two patch masses, two bounded rail
factors, and a port vector. It is bounded by
$C\Gamma_L(l)^{-2j}g_L(l)^{-2}\leq C\Gamma_L(l)^{-2}$. Thus the sum is
uniformly bounded, since
\[
 \sum_{l\in\Lambda_L}\Gamma_L(l)^{-2}
 \leq\left(\sum_{n\in\mathbb Z}(1+|n|)^{-2}\right)^4\leq81.
\]
The extra constant $k^{-1}$ is finite. The divisor, scaled signals, and
all zero-alias ratios in the correction and notch are bounded by the
preceding train lemma. Their formulas contain only finitely many sums and
products of these quantities, so the triangle inequality and the
finite-dimensional matrix product bound complete the uniform estimate.
\end{proof}

\subsection{Closing the two wire branches}

\begin{proposition}[Entrywise weighted wire bound]
\label{prop:ex-lattice-wire-entrywise-bound}
For fixed stages $j,t\geq4$ and $k>0$, each scalar entry of the dual
stencil is holomorphic on one common complex cube and satisfies
\[
 |W_{P,0}(p)|\leq C_0,\qquad
 |W_{P,l}(p)|\leq C_1\Gamma_L(l)^{-j}\quad(l\ne0),
\]
uniformly in odd $L\geq3$. Hence the unified bound is
$C\Gamma_L(l)^{-j}$.
\end{proposition}

\begin{proof}
Use the zero and nonzero branch formulas in
Section~\ref{sec:ex-lattice-free-kernel-estimates}. In the zero branch,
Lemma~\ref{lem:ex-lattice-mass-flux-weight} and
Lemma~\ref{lem:ex-lattice-tail-notch-weights} show that every factor has
weight one. In the first nonzero branch product, the scaled signal,
port vector, and flux contribute respectively
$g_L(l)$, $g_L(l)^{-2}$, and $\Gamma_L(l)^{-j}$. Their product is
at most $\Gamma_L(l)^{-j}$. The second nonzero product has levy weight
one and flux weight $\Gamma_L(l)^{-j}$. Finite intermediate matrix sums
multiply the constant by only the fixed dimension of the face spaces.
The branch constants and widths can be replaced by their maximum and
minimum, respectively. This proves the assertion.
\end{proof}
